\chapter{ساخت و درک ربات}

\begin{goalbox}
توکن را به تله‌بات متصل کنید و بفهمید هر پیام چگونه پاسخ خود را می‌گیرد.
\end{goalbox}

فایل کامل از قبل در مخزن قرار دارد. خواندن آن از بالا به پایین، سریع‌ترین راه
برای ارتباط‌دادن هر مفهوم با کد عملی است.

\begin{LTR}
\lstinputlisting[language=Python,caption={The complete code in code/bot.py}]{\CodePath/bot.py}
\end{LTR}

\section{گام اول: واردکردن ابزارها}

ماژول
\lr{os}
متغیر محیطی را می‌خواند.
\lr{telebot}
برنامه‌ای می‌سازد که با تلگرام صحبت می‌کند و
\lr{types}
نوع پیام‌ها را برای خواننده و ویرایشگر مشخص می‌کند. این نوع‌ها رفتار ربات را
تغییر نمی‌دهند.

\section{گام دوم: پاسخ‌هایی که آسان آزموده می‌شوند}

تابع‌های
\lr{start\_reply}
و
\lr{help\_reply}
و
\lr{echo\_reply}
فقط متن می‌سازند و به اینترنت نیاز ندارند. جداکردن آن‌ها گفت‌وگوی مورد انتظار
را روشن می‌کند و به آزمون‌های آفلاین اجازه می‌دهد نتیجه را بررسی کنند.

\section{گام سوم: ثبت پاسخ‌دهنده‌ها}

پاسخ‌دهنده یا
\lr{handler}
فقط تابعی است که برای یک نوع پیام ورودی انتخاب می‌شود. دستورهای بالای تابع‌ها
این انتخاب را مشخص می‌کنند:

\begin{itemize}
\item عبارت
\lr{commands=["start"]}
دستور
\lr{/start}
را مدیریت می‌کند؛
\item عبارت
\lr{commands=["help"]}
دستور
\lr{/help}
را مدیریت می‌کند؛
\item عبارت
\lr{content\_types=["text"]}
متن‌های دیگر را مدیریت می‌کند؛
\item آخرین پاسخ‌دهنده، پیام‌های رایج غیرمتنی را پوشش می‌دهد.
\end{itemize}

ترتیب مهم است: پاسخ‌دهنده‌های دستور پیش از پاسخ‌دهندهٔ عمومی متن قرار گرفته‌اند؛
پس هر دستور پاسخ ویژهٔ خودش را می‌گیرد.

\section{گام چهارم: آغاز دریافت پیام}

تابع
\lr{main}
در صورت نبودن
\lr{BOT\_TOKEN}
با خطایی خوانا متوقف می‌شود. اگر توکن وجود داشته باشد، تابع
\lr{infinity\_polling}
تا زمان فشردن
\lr{Ctrl+C}
به دریافت پیام ادامه می‌دهد. گزینهٔ
\lr{skip\_pending=True}
پیام‌های قدیمی جمع‌شده در زمان خاموش‌بودن این ربات آموزشی را نادیده می‌گیرد.

\begin{checkpointbox}
اکنون می‌توانید پاسخ‌دهندهٔ دقیق دستورهای
\lr{/start}
و
\lr{/help}
و همچنین متن عادی و عکس را در کد نشان دهید.
\end{checkpointbox}
